%%%%%%%%%%%%%%%%%%%%%%%%%%%%%%%%%%%%%%%%%%%%%%%%%%%%%%%%%%%%%%%%%%%
% Weekly Report: created by MW, 2003 and Modified by Joohyun, 2009
%%%%%%%%%%%%%%%%%%%%%%%%%%%%%%%%%%%%%%%%%%%%%%%%%%%%%%%%%%%%%%%%%%%
\documentclass[10pt,letterpaper,oneside]{article}
\usepackage{cite,graphicx,subfigure,amsmath,amsthm,amssymb,epsfig,pifont}
\usepackage[UTF8]{ctex} % 加载 ctex 宏包

\setlength{\textwidth}{15.25cm}%
\setlength{\textheight}{23cm}%
\setlength{\oddsidemargin}{0 cm}%
\setlength{\evensidemargin}{0 cm}%
\setlength{\topskip}{0 cm}%
\setlength{\topmargin}{-1 cm}%
\setlength{\headheight}{0 cm}%

%\usepackage{eufrak}
% \pagestyle{empty}

\textwidth  6.6in
\textheight 9.1in
\setlength{\oddsidemargin}{-0.04 in}
\setlength{\topmargin}{-0.7in}

\newcommand{\weekly}[5]{
    \pagestyle{plain}
    \newpage
    \setcounter{page}{1}
    \noindent
    \begin{center}
    \begin{tabular}[t]{||c||c||c||}
    \hline\hline\hspace{3mm}\rule[-3mm]{0mm}{10mm}
    \textbf{\Large{#3 #2}}\hspace{3mm} &\hspace{3mm}\textbf{\Large{Report\##5 - #4/#2}}\hspace{3mm} &\hspace{3mm}\textbf{\Large{#1}}\hspace{3mm}
    \\ \hline\hline
    \end{tabular}
    \end{center}
    %\markboth{#1}{#1}
    \vspace*{4mm}
}

\begin{document}

%% Main Title Section %%
% \weekly{Wei Guanzhao}{2026}{Summer}{07/11}{04}
\weekly{Wei Guanzhao}{2026}{Summer}{07/17}{05}




%%%%%%%%%%%%%%%% I. Task achieved in this week section %%%%%%%%%%%%%%%%%
\vspace{1 mm}
\begin{picture}(10,0)\linethickness{0.5 mm}
\put(-25,0){\line(1,0){500}}
\end{picture}
\vspace{2 mm}

\noindent\textbf{\large{I. Tasks achieved Last Week}}

\begin{picture}(10,0) \linethickness{0.5 mm}
\put(-25,0){\line(1,0){500}}
\end{picture}
%%%%%%%%%%%%%%%%%%%%%%%%%%%%%%%%%%%%%%%%%%%%%%%%%%%%%%%%%%%%%%%%%%%%%%
% \begin{itemize} %\setlength\itemsep{-\parsep}   %
%   \item[\normalsize{\ding{110}}] Fixed bugs in the AerialVLN code and started training.  \vspace{-3 mm}\\
%   % \subitem {\footnotesize{\ding{108}}} 
%   % \item[\normalsize{\ding{110}}] 跟蒋老师学习具身智能  \vspace{-3 mm}\\
% \end{itemize}                      % --> To Here
\begin{itemize} %\setlength\itemsep{-\parsep}   %
  \item[\normalsize{\ding{110}}] Think about how to revise and extend the paper for journal submission  \vspace{-3 mm}\\
  % \subitem {\footnotesize{\ding{108}}} 
  \item[\normalsize{\ding{110}}] Survey the development of World model in UAV vision-language navigation  \vspace{-3 mm}
\end{itemize}                      % --> To Here
%%%%%%%%%%%%%%%%%%%%%%%%%%%%%%%%%%%%%%%%%%%%%%%%%%%%%%%%%%%%%%%%%%%%%%



%%%%%%%%%%%%%%%%% II. Feedback and Interaction section %%%%%%%%%%%%%%%
% \vspace{1 mm}
% \begin{picture}(10,0)\linethickness{0.5 mm}
% \put(-25,0){\line(1,0){500}}
% \end{picture}
% \vspace{2 mm}

% \noindent \textbf{\large{II. Feedback and Interaction}}

% \begin{picture}(10,0) \linethickness{0.5 mm}
% \put(-25,0){\line(1,0){500}}
% \end{picture}
% %%%%%%%%%%%%%%%%%%%%%%%%%%%%%%%%%%%%%%%%%%%%%%%%%%%%%%%%%%%%%%%%%%%%%%

% \begin{itemize}\setlength\itemsep{-\parsep}
%     \item[\normalsize{\ding{110}}] \textbf{\emph{Prof. Kuo's F/I}}
%     \setlength\itemsep{-\parsep}\vspace{-2mm}
%     \begin{itemize} \setlength\itemsep{-\parsep}
%         \item  \emph{See the first and last paragraph in the weekly report}\vspace{-3mm}\\
%     \end{itemize}

% \end{itemize}
%%%%%%%%%%%%%%%%%%%%%%%%%%%%%%%%%%%%%%%%%%%%%%%%%%%%%%%%%%%%%%%%%%%%%%

%%%%%%%%%%%%%%%%%%%%% III. Report Section %%%%%%%%%%%%%%%%%%%%%%%%%%%%
\vspace{1 mm}
\begin{picture}(10,0)\linethickness{0.5 mm}
\put(-25,0){\line(1,0){500}}
\end{picture}
\vspace{2 mm}

\noindent \textbf{\large{II. Report}}

\begin{picture}(10,0) \linethickness{0.5 mm}
\put(-25,0){\line(1,0){500}}
\end{picture}

%%%%%%%%%%%%%%%%%%%%%%%%%%%%%%%%%%%%%%%%%%%%%%%%%%%%%%%%%%%%%%%%%%%%%%
\begin{enumerate}\setlength\itemsep{-\parsep}
  % \item \Large\textbf{数据收集 \& 修改训练代码} 
  \item \Large\textbf{Revise and extend the paper for journal submission} 
  
  % \large\textbf{Data Collection:} 
  \normalsize{
      Analyze the current methodological issues in the paper, consider a V2 version, 
      find a suitable journal, add results on the AerialVLN dataset, and submit.
  }

  % \large\textbf{Training Code Modification:} 
  \normalsize{
    % (Commented out original Chinese content, kept as is)
  }

  \item \Large\textbf{Survey the development of world models in UAV vision-language navigation}
  
  \normalsize{
    The core positioning of world models is to shift from "reactive execution" to "predictive planning".
    \begin{itemize}
      \item Explicitly transform abstract natural language instructions into future first-person visual sequences;
      \item Provide motion priors and spatial foresight. In complex urban scenes (severe occlusion, drastic viewpoint changes), 
      world model allow the agent to "see" the environment around turns in advance by generating future frames, 
      avoiding the "short-sightedness" of pure historical observation policies.
      \item Leverage imagined future dynamic features to optimize decision-making, 
      while world model adjust generated content according to planned trajectories, achieving closed-loop collaborative training.
    \end{itemize} 

    There are also some issues to be addressed:
    \begin{itemize}
      \item Significant inference latency and computational burden (core pain point);
      \item Insufficient pixel-level generation quality and spatiotemporal consistency;
      \item Sim-to-Real transfer robustness;
    \end{itemize} 
  }
    % \begin{itemize}
    %   \item 原始的指令文本
    %   \item RGB 图像 
    % \end{itemize} 
  % 原始收集的数据集中，只保存了处理好的 RGB 图像特征，未保留原始 RGB 图像，还有原始指令文本也需要保留，训练时需要使用大模型Tokenizer 进行处理。

  % \item \Large\textbf{跟蒋老师学习具身智能}
  % \normalsize{
  %   学习了应该如何去定义具身智能。对于具身智能来说，大模型相当于大脑，
  %   大脑是相当重要的，但是还要把身体包含进来，身体是具身智能最重要的
  %   一部分，因为具身智能需要与环境进行感知，所感知的信息传输给大脑
  %   进行信息的提取和转换成能够理解的知识。此外，如何能够让智能体
  %   在与环境交互的过程中持续的学习知识，而不是利用预先定义好的策略
  %   让智能体重复的做同样的事情也是未来需要关注的事情。Sim2Real 相关的
  %   实验具有一定的局限性，对于算力和数据的高精度性要求较高，且对于智能体
  %   的泛化性有一定的挑战，稍微偏移的场景可能会产生不可控制的结果。
  %   对于未来的想法，强化学习和模仿学习是可以值得探索的空间。
  %   最后，对于我们想要做的算法，建议要上真机实验，模拟器上的空间和物体
  %   特权弱化了真实的场景，对于真实的机器人来说，底层控制和上层算法同样
  %   重要。
  % }

  % \item \Large\textbf{无人机视觉语言导航}
  \normalsize{
    % 无人机视觉语言导航主要挑战有：跨模态语义对齐困难、
    % 三维空间推理能力不足、
    % 长时序空间推理与累积误差、机载算力严重受限等。现在大多
    % 方法主要集中在解决长时序空间推理与累积误差这个问题，为
    % 解决这个问题，应该提出一个有效整合时空上下文信息的方法。
    % 目前模型是一次性接收指令后“开环”执行。如果可以引入
    % “闭环”交互，让无人机在导航过程中能主动提问来消除歧义。
    % 这种交互过程本身就会产生新的时空上下文。这些新信息会
    % 被用于更新记忆并指导后续导航，形成一个
    % “感知-行动-对话-再感知”的增强回路。
  }


 
  % \item \large\textbf{Literature Review} \\
  % \normalsize{Contents}\\
  % Title...\\
  % Content....\\
  % Content....\\

%%%%%%%%%%%%%%%%%%%%%%%%%%%%%%%%%%%%%%%%%%%%%%%%%%%%%%%%%%%%%%%
%\begin{figure}[ht]
%\centering \scalebox{0.6}[0.6]{\includegraphics{fig/etx3}}
%\caption{Expected transmission counts
%}\label{fig:maxduNei}\vspace{-1 mm}
%\end{figure}
%%%%%%%%%%%%%%%%%%%%%%%%%%%%%%%%%%%%%%%%%%%%%%%%%%%%%%%%%%%%%%%

% \item \textbf{(Report B)} Journal 1 final result\\
% \begin{itemize} \setlength\itemsep{-\parsep}\vspace{-5mm}
%     \item Blablabla\cite{ts.36.104} shows blabla \vspace{-1mm}\\
%     \item Based on the result of \cite{rhpark05}, we can blabla  \vspace{-1mm}\\
%     \item In contrast with the algorithm introduced by \cite{Sheldon&Seo&Torkildson&Rodwell09}, Hit and Crook Algorithm (HCA) we present here has blablas   \vspace{-1mm}\\
% \end{itemize}

%%%%% If you need to add references related to report, use this %%
\bibliographystyle{IEEEtran}
%\bibliography{IEEEabrv,johndoe}
\bibliography{IEEEabrv,836673citation}

\vspace{3 mm}

\end{enumerate}


%%%%%%%%%%%%%%%%% IV. Plan for the next week section %%%%%%%%%%%%%%%%%%%%
\vspace{1 mm}
\begin{picture}(10,0)\linethickness{0.5 mm}
\put(-25,0){\line(1,0){500}}
\end{picture}
\vspace{2 mm}

\noindent\textbf{\large{III. Plan for the next week}}

\begin{picture}(10,0) \linethickness{0.5 mm}
\put(-25,0){\line(1,0){500}}
\end{picture}
%%%%%%%%%%%%%%%%%%%%%%%%%%%%%%%%%%%%%%%%%%%%%%%%%%%%%%%%%%%%%%%%%%%%%%
%\begin{itemize}\setlength\itemsep{-\parsep}   %
%\begin{itemize}\setlength\itemsep{-\parsep}\vspace{-2mm}
% \begin{itemize}
%   \item[\normalsize{\ding{110}}] The bugs have been fixed. 
%   First, run 1-2 epochs on the server to check whether 
%   there are any other issues. 
%   Then rent an 8-GPU server for training and hyperparameter tuning. \vspace{-3 mm}
%   \item[\normalsize{\ding{110}}] Read CVPR
%    papers and think about the specific implementation process of closed-loop interactive navigation. \vspace{-3 mm}\\
%   % \subitem {\footnotesize{\ding{108}}} I believe you understand what I mean\\
 
% \end{itemize}                      % --> To Here

\begin{itemize}
  \item[\normalsize{\ding{110}}] 
    Survey potential journals for extension, and think about how to improve the methodological issues of the paper \vspace{-3 mm}\\
  \item[\normalsize{\ding{110}}] 
    Read papers and think about problems in UAV visual dialogue navigation.
    \begin{itemize}
      \item Study improving the timing of question asking
      \item Study the structure of model-generated questions
      \item Study memory mechanisms based on topological graphs or other data structures
    \end{itemize}
    % \vspace{-3 mm}\\
  % \subitem {\footnotesize{\ding{108}}} I believe you understand what I mean\\
\end{itemize}       
%%%%%%%%%%%%%%%%%%%%%%%%%%%%%%%%%%%%%%%%%%%%%%%%%%%%%%%%%%%%%%%%%%%%%%


%% Milestone Section %%
% \vspace{1 mm}
% \begin{picture}(10,0)\linethickness{0.5 mm}
% \put(-25,0){\line(1,0){500}}
% \end{picture}
% \vspace{2 mm}

% \noindent \textbf{\large{V. Milestone}}

% \begin{picture}(10,0) \linethickness{0.5 mm}
% \put(-25,0){\line(1,0){500}}
% \end{picture}
% %%%%%%%%%%%%%%%%%%%%%%%%%%%%%%%%%%%%%%%%%%%%%%%%%%%%%%%%%%%%%%%%%%%%%%
% \\ \\
% \textbf{\emph{Main Goal of this Semester : Publish a paper to Whiskey conference chaired by Dr. Quack} }\\
% Subgoal : Keeping motivation \\
% \\ \textbf{\emph{Plan to achieve the goal}}


% \begin{itemize} \setlength\itemsep{-\parsep}% <-- Modify the contents  for 'V. Milestone'

% \item ...

% \item ...
% %    \subitem -
% \end{itemize}
%%%%%%%%%%%%%%%%%%%%%%%%%%%%%%%%%%%%%%%%%%%%%%%%%%%%%%%%%%%%%%%%%%%%%%

%\newpage
%
%\begin{picture}(10,0) \linethickness{0.5 mm}
%\put(-25,0){\line(1,0){450}}
%\end{picture}
%\begin{center}
%\textbf{\large{Appendix}} ~~Aging analysis of topology
%construction schemes
%\end{center}
%
%\begin{picture}(10,0) \linethickness{0.5 mm}
%\put(-25,0){\line(1,0){450}}
%\end{picture}
%\vspace{3mm}\\


\end{document}
